\documentclass[a4paper,12pt]{article}
\usepackage[utf8]{inputenc}
\usepackage[T1]{fontenc}
\usepackage[ngerman]{babel}
\usepackage{amsmath}
\usepackage{amssymb}
\usepackage{enumitem}
\usepackage{geometry}
\geometry{a4paper, margin=2.5cm}

\title{Einsendeaufgaben -- Aufgabe 1.1\\[0.5em]
\large 61211 Analysis}
\author{}
\date{}

\begin{document}
\maketitle

\section*{Aufgabe 1}

\begin{enumerate}[label=\alph*)]

    \item
    \begin{enumerate}[label=(\roman*)]

        \item
        \underline{1. $\forall x \in \mathbb{R}^3 : \|x\| \geq 0$}

        Da $\|\cdot\|_2$ und $|\cdot|$ Normen sind, gilt
        \[
            \|(x_1, x_3)\|_2 \geq 0 \quad \text{und} \quad |x_2| \geq 0.
        \]
        Also
        \[
            \|x\| = \|(x_1, x_3)\|_2 + |x_2| \geq 0.
        \]

        \underline{2. $\forall x \in \mathbb{R}^3 : \|x\| = 0 \Longleftrightarrow x = 0$}

        \glqq$\Rightarrow$\grqq-Richtung:

        Sei $\|x\| = 0$. Dann gilt
        \[
            \|x\| = \|(x_1, x_3)\|_2 + |x_2| = 0.
        \]
        Da $\|(x_1, x_3)\|_2 \geq 0$ und $|x_2| \geq 0$, folgt
        \[
            \|(x_1, x_3)\|_2 = 0 \quad \text{und} \quad |x_2| = 0.
        \]
        Da $\|\cdot\|_2$ und $|\cdot|$ Normen sind, folgt
        $(x_1, x_3) = 0$ und $x_2 = 0$,
        also $x = \begin{pmatrix}x_1\\x_2\\x_3\end{pmatrix} = 0$.

        \glqq$\Leftarrow$\grqq-Richtung:

        Sei $x = 0$. Da $\|\cdot\|_2$ und $|\cdot|$ Normen sind, folgt
        $\|(x_1, x_3)\|_2 = 0$ und $|x_2| = 0$. Demnach gilt
        \[
            0 = \|(x_1, x_3)\|_2 + |x_2| = \|(x_1, x_2, x_3)\|.
        \]

        \item
        $\forall x \in \mathbb{R}^3\; \forall a \in \mathbb{R} : \|ax\| = |a|\,\|x\|$
        \begin{align*}
            \|ax\| &= \|(ax_1, ax_3)\|_2 + |ax_2| \\
                   &= \|a(x_1, x_3)\|_2 + |ax_2| \\
                   &= |a|\,\|(x_1, x_3)\|_2 + |a|\,|x_2| \\
                   &= |a|\bigl(\|(x_1, x_3)\|_2 + |x_2|\bigr) \\
                   &= |a|\,\|x\|
        \end{align*}

        \item
        $\forall x, y \in \mathbb{R}^3 : \|x + y\| \leq \|x\| + \|y\|$
        \begin{align*}
            \|x + y\| &= \|(x_1 + y_1, x_3 + y_3)\|_2 + |x_2 + y_2| \\
                      &\leq \bigl(\|(x_1, x_3)\|_2 + \|(y_1, y_3)\|_2\bigr)
                            + \bigl(|x_2| + |y_2|\bigr) \\
                      &= \bigl(\|(x_1, x_3)\|_2 + |x_2|\bigr)
                         + \bigl(\|(y_1, y_3)\|_2 + |y_2|\bigr) \\
                      &= \|x\| + \|y\|
        \end{align*}

    \end{enumerate}

    \item
    \begin{enumerate}[label=(\roman*)]

        \item
        Es gilt $\sqrt{x_1^2 + x_3^2} \leq \sqrt{x_1^2 + x_2^2 + x_3^2}$
        und $|x_2| \leq \sqrt{x_1^2 + x_2^2 + x_3^2}$, da:

        \begin{enumerate}[label=\arabic*.]
            \item
            $\sqrt{x_1^2 + x_3^2} \leq \sqrt{x_1^2 + x_2^2 + x_3^2}$
            \[
                \Longleftrightarrow x_1^2 + x_3^2 \leq x_1^2 + x_2^2 + x_3^2
                \Longleftrightarrow 0 \leq x_2^2
            \]

            \item
            $|x_2| \leq \sqrt{x_1^2 + x_2^2 + x_3^2}$
            \[
                \Longleftrightarrow x_2^2 \leq x_1^2 + x_2^2 + x_3^2
                \Longleftrightarrow 0 \leq x_1^2 + x_3^2
            \]
        \end{enumerate}

        Wir wählen $c_1 = 2$ und daraus folgt:
        \begin{align*}
            \|(x_1, x_2, x_3)\| &= \|(x_1, x_3)\|_2 + |x_2| \\
                                 &\leq \|(x_1, x_2, x_3)\|_2 + \|(x_1, x_2, x_3)\|_2 \\
                                 &= 2\,\|(x_1, x_2, x_3)\|_2
        \end{align*}

        \item
        Wir setzen $c_2 = 3$ und daraus folgt:

        Es gilt $\sqrt{x_1^2 + x_3^2} \leq 2\,\|(x_1, x_2, x_3)\|_\infty$
        und $|x_2| \leq \|(x_1, x_2, x_3)\|_\infty$, da:

        \begin{enumerate}[label=\arabic*.]
            \item
            \[
                x_1^2 + x_3^2 \leq 2\max\{x_1^2, x_3^2\}
                \leq 2\max\{x_1^2, x_2^2, x_3^2\}
                = 2\,\|(x_1, x_2, x_3)\|_\infty^2
            \]
            Demnach gilt $\sqrt{x_1^2 + x_3^2} \leq 2\,\|(x_1, x_2, x_3)\|_\infty$.

            \item
            \[
                |x_2| \leq \max\{|x_1|, |x_2|, |x_3|\} = \|(x_1, x_2, x_3)\|_\infty
            \]
        \end{enumerate}

        Daraus folgt:
        \begin{align*}
            \|(x_1, x_2, x_3)\| &= \sqrt{x_1^2 + x_3^2} + |x_2| \\
                                 &\leq 2\,\|(x_1, x_2, x_3)\|_\infty
                                       + \|(x_1, x_2, x_3)\|_\infty \\
                                 &= 3\,\|(x_1, x_2, x_3)\|_\infty
        \end{align*}

    \end{enumerate}

\end{enumerate}

\end{document}
